\documentclass[12pt]{article}
\usepackage[utf8]{inputenc}
\usepackage[T1]{fontenc}
\usepackage[margin=1in]{geometry}
\usepackage{setspace}
\usepackage{xcolor}
\usepackage[colorlinks=true, urlcolor=blue!60!black, linkcolor=blue!60!black]{hyperref}
\setlength{\parindent}{2em}
\setlength{\parskip}{1.25em}
\renewcommand{\baselinestretch}{1.0}
\begin{document}
\begin{center}
{\LARGE \textit{Dames \& Moore v. Regan}, 453 U.S. 654 (1981)\par}
\vspace{0.8em}
{\large Case Brief --- Four-Page Edition\par}
\end{center}
\vspace{1em}

\section{Memory Jogger}

During the Iran Hostage Crisis, the Court upheld executive orders nullifying attachments on frozen Iranian assets under IEEPA and suspending private U.S. lawsuits in favor of an international claims tribunal, resting the suspension on congressional acquiescence and the tenor of related statutes (\href{https://supreme.justia.com/cases/federal/us/453/654/}{Justia}).

\section{Detailed Case Facts}

After Iranian militants seized the U.S. Embassy in Tehran on November 4, 1979, President Carter declared a national emergency under the International Emergency Economic Powers Act (IEEPA) and blocked all Iranian government property subject to U.S. jurisdiction. Treasury regulations declared any unlicensed attachment of blocked property null and void and provided that any license could be ``amended, modified, or revoked at any time,'' then issued general licenses permitting prejudgment attachments. Dames \& Moore, an engineering firm owed some \$3.4 million for nuclear-plant site studies, sued Iran and obtained prejudgment attachments against Iranian banks.

On January 19, 1981, the United States and Iran concluded the Algiers Accords to free the 52 hostages. The United States agreed to terminate legal proceedings by U.S. nationals against Iran, nullify attachments, transfer blocked assets, and submit claims to a binding Iran--United States Claims Tribunal backed by a \$1 billion security account. Carter's executive orders revoked the attachment licenses and ordered assets transferred; President Reagan ratified them and ``suspending'' all claims eligible for the Tribunal, which would revive if the Tribunal lacked jurisdiction.

\section{Procedural History}

Dames \& Moore won summary judgment against Iran, but the district court stayed execution and vacated the attachments in light of the executive orders. Dames \& Moore then sued the United States and Treasury Secretary Regan to block the orders; the district court dismissed. Facing the Accords' July 19, 1981 transfer deadline, the Court granted certiorari before judgment, heard argument on June 24, and decided the case on July 2, 1981.

\section{Judicial Votes}

Affirmed, unanimous on the core executive-power judgment.
\begin{itemize}
\item \textbf{Majority:} Justice Rehnquist, joined by Chief Justice Burger and Justices Brennan, Stewart, White, Marshall, and Blackmun; Stevens joined all but Part V; Powell joined all but footnote 6.
\item \textbf{Concurring in part:} Justice Stevens (would not reach Tucker Act jurisdiction).
\item \textbf{Concurring in part and dissenting in part:} Justice Powell (dissenting from footnote 6's no-taking conclusion).
\end{itemize}

\section{Holding}

IEEPA expressly authorizes the President, in a declared emergency, to nullify post-freeze attachments and compel transfer of frozen foreign assets---\textit{Youngstown} Category One. Although no statute specifically authorizes suspending private suits, where claims settlement is a necessary incident to resolving a major foreign-policy dispute, an alternative forum exists, and Congress has long acquiesced in executive claims settlement, the President may suspend pending claims without divesting Article III jurisdiction. A takings suit under the Tucker Act remains available in the Court of Claims. The Court disclaimed any ``plenary power'' to settle all claims.

\section{Analysis of Opinions}

\textbf{Majority (Rehnquist).} The Court adopted Justice Jackson's \textit{Youngstown} framework but recast its three categories as a ``spectrum running from explicit congressional authorization to explicit congressional prohibition.'' \emph{Attachments:} IEEPA's text empowers the President to ``direct and compel, nullify, void, prevent or prohibit'' transfers of foreign property; frozen assets are a diplomatic ``bargaining chip'' that private litigants could not dismantle through revocable attachments, so the orders enjoyed the ``strongest of presumptions.'' Footnote 6 held that attachments under revocable licenses conferred no compensable property interest. \emph{Suspension:} neither IEEPA nor the Hostage Act authorized suspending suits, but their ``general tenor,'' an unbroken history of executive claims settlement back to 1799, and the International Claims Settlement Act showed congressional acquiescence; under Frankfurter's \textit{Youngstown} ``gloss'' and \textit{Midwest Oil}, long practice raised a presumption of consent. Suspension changed the substantive rule of law rather than stripping jurisdiction, and the Tribunal offered a meaningful alternative forum. \emph{Part V:} the takings question was unripe, but Tucker Act jurisdiction was available.

\textbf{Stevens (concurring in part).} The chance of a taking was ``so remote that I would not address the jurisdictional question considered in Part V.''

\textbf{Powell (concurring in part and dissenting in part).} Under \textit{Armstrong v. United States}, an attachment is a recognized property right, and it was premature to foreclose takings claims in hundreds of pending cases. Foreign-affairs powers cannot displace the Just Compensation Clause, which bars forcing a few citizens to bear burdens that belong to the public as a whole.

\section{Examples: Future Applications}

\begin{enumerate}
\item \textbf{Same-Side: Executive Agreement Preempting State Law.} A sole executive agreement creating a compensation foundation for historical claims preempts a conflicting state insurance-disclosure law, as \textit{American Insurance Ass'n v. Garamendi} later held.
\item \textbf{Same-Side: Revoking Licensed Attachments.} Under an IEEPA sanctions program, the President revokes a license that let contractors attach frozen accounts and transfers the funds to an escrow for a peace accord: Category One, valid.
\item \textbf{Opposite-Side: Directive to State Courts.} A presidential memorandum ordering state courts to disregard procedural-default rules to enforce a non-self-executing treaty has no ``systematic, unbroken, executive practice'' behind it and exceeds Article II, as \textit{Medell\'in v. Texas} held.
\item \textbf{Opposite-Side: Defying a Statutory Bar.} If Congress expressly forbids transferring a terror sponsor's frozen assets, a contrary executive order sits at Category Three, where acquiescence cannot help.
\item \textbf{Fence-Sitter: IEEPA Tariffs.} Sweeping emergency tariffs under IEEPA: supporters invoke \textit{Dames \& Moore}'s broad reading of IEEPA; opponents answer that tariffs are a core Article I power lacking the historical gloss of claims settlement and triggering major-questions scrutiny.
\end{enumerate}

\section{Critique}

\textbf{Silence as authorization.} Having found that no statute authorized suspending suits, the Court used their ``general tenor'' to supply authority. Critics note that IEEPA was enacted in 1977 to \textit{curtail} emergency powers, and that treating congressional inaction as consent ignores bicameralism and the practical impossibility of overriding a President in a hostage crisis.

\textbf{Structure and rights.} Progressive critics say the decision aggrandized the presidency and echo Powell's concern for private property and remedies; formalists say overriding pending causes of action should require a treaty or statute, and warn that historical gloss creates a one-way ratchet favoring Article II---a concern \textit{Medell\'in} later answered by confining the case to claims settlement.

\section{Key Quotations}

\begin{quote}
``Past practice does not, by itself, create power, but `long-continued practice, known to and acquiesced in by Congress, would raise a presumption that the [action] had been [taken] in pursuance of its consent. . . .'\,''
\end{quote}
\noindent --- Justice Rehnquist, for the Court, 453 U.S. at 686.

\begin{quote}
``Finally, we reemphasize the narrowness of our decision. We do not decide that the President possesses plenary power to settle claims, even as against foreign governmental entities.''
\end{quote}
\noindent --- Justice Rehnquist, for the Court, 453 U.S. at 688.

\begin{quote}
``The Government must pay just compensation when it furthers the Nation's foreign policy goals by using as `bargaining chips' claims lawfully held by a relatively few persons and subject to the jurisdiction of our courts. The extraordinary powers of the President and Congress upon which our decision rests cannot, in the circumstances of this case, displace the Just Compensation Clause of the Constitution.''
\end{quote}
\noindent --- Justice Powell, concurring in part and dissenting in part, 453 U.S. at 691.

\end{document}
